\documentclass[12pt]{article}

% ── Packages ──────────────────────────────────────────────────────────────────
\usepackage[margin=1in]{geometry}
\usepackage{setspace}
\usepackage{titlesec}
\usepackage{booktabs}
\usepackage{graphicx}
\usepackage{amsmath}
\usepackage{hyperref}
\usepackage{parskip}
\usepackage{xcolor}

\hypersetup{
    colorlinks=true,
    linkcolor=black,
    urlcolor=blue,
}

% ── Title Page Info ───────────────────────────────────────────────────────────
\title{
    \vspace{2cm}
    {\Large AML S26 Final Project} \\[0.5em]
    {\large Jordan Miller and Sydney Williams}
}
\author{}
\date{May 2026}

% ── Placeholder style ─────────────────────────────────────────────────────────
\newcommand{\placeholder}[1]{\textit{\textcolor{gray}{[#1]}}}

% =============================================================================
\begin{document}
% =============================================================================

\maketitle
\thispagestyle{empty}
\newpage

\tableofcontents
\newpage


% =============================================================================
\section{Part I}
% =============================================================================

\subsection{Data Overview}

\subsubsection{File size, class balance, scale of features}

\placeholder{File size, class balance, scale of features}


\subsection{Design Decisions}

\subsubsection{Pipeline elements}

\placeholder{Talk about why you picked the elements of the pipeline you used}

\subsubsection{Model selection}

\placeholder{Talk about what models picked and why}


\subsection{Traditional Results}

\subsubsection{Development and training time}

\placeholder{Talk about dev and training time}

\subsubsection{Numerical metrics}

\placeholder{Talk about numerical metrics (F1, roc auc, etc)}

\subsubsection{Confusion matrix}

\placeholder{Talk about confusion matrix}

\subsubsection{Predicted results on test set}

\placeholder{Talk about predicted results on test set}


\subsection{Deep Results}

\subsubsection{Development and training time}

\placeholder{Talk about dev and training time}

\subsubsection{Numerical metrics}

\placeholder{Talk about numerical metrics (F1, roc auc, etc)}

\subsubsection{Confusion matrix}

\placeholder{Talk about confusion matrix}

\subsubsection{Predicted results on test set}

\placeholder{Talk about predicted results on test set}


\subsection{Conclusion}

\subsubsection{Model comparison}

\placeholder{Compare two models}

\subsubsection{Lessons learned}

\placeholder{What have you learned?}


% =============================================================================
\section{Part II}
% =============================================================================

\subsection{Data Overview}

\subsubsection{File size, class balance, scale of features}

The PartII\_dev.csv file contains 125 features. Each feature has 158 numerical samples, with 126 belonging to class 0 and 32 belonging to class 1. This represents a 4:1 class imbalance. The features have widely ranging scales. There are no missing values.


\subsection{Design Decisions}

\subsubsection{Pipeline elements}

The massive redundancy in features and very small sample size were the main constraints for this part of the project.

SMOTE was used for both the traditional and deep learning processes. SMOTE (Synthetic Minority Over-Sampling Technique)synthesizes new minority-class samples by interpolating between existing ones in the feature space. With only 32 class 1 samples, a classifier that always predicts class 0 would have 80\% accuracy. Therefore, SMOTE is necessary to give the model a better view of the boundary between classes and prevent it from only predicting class 0. 

For the traditional model, StratifiedKFold with 5 folds was used to compare models. It compared performance of SVM with PCA, LogReg with PCA, RandomForest, and HistGBM.

Once the best model was selected, GridSearchCV was used to tune model parameters.

Due to the large number of features and wide variation in scale, component scaling and reduction was used for the traditional portion. First, StandardScaler was used to scale the features, then PCA was used to reduce the number of components to those that represent 95\% of the variance (11 features).

For the deep learning section, BatchNorm was used after each layer. BatchNorm normalizes activations to the zero mean and unit variance per batch, which acts a regularizer. It also reduces sensitivity to weight activation.

ReLU activation was used, which is the standard choice. Initially, dropout was set to 0.5. This meant that after each BatchNorm+ReLU a random 50\% selection of neurons would be zeroed. It was chosen to have such aggressive regularization because of the small dataset.

\subsubsection{Model selection}

Initially, SVM with PCA was chosen for the traditional method and MLP (Multi-Layer Perceptron)was chosen for the deep learning method.

SVM (Support Vector Machine) was chosen based on the results of the model comparison CV. SVMs work well for small-sample, high-dimensional problems. The RBF kernel was chosen for it's ability to map into an infinite-dimensional space. This allows non-linear decision boundaries, which is appropriate for this problem as we have no knowledge of whether the classes are really linearly separable.

PCA was chosen to decorrelate the features and reduce dimensionality as outlined in the section above. It also plays very nicely with SVMs. Redundant features can make SVM optimization unstable and greatly reduces the training speed of the model. PCA discarding the dimensions that contain the remaining 5\% of variance can also improve test performance if those dimensions are noisy.

For the deep learning portion, MLP was chosen because it typically does better than other deep learning models when learning non-linear interactions between features. The dataset for part two is very small, so it is anticipated to perform worse for deep learning compared to traditional learning in general.

It was decided to have cap the model at three hidden layers, starting with a width of 128 and halving the width each layer (128 -> 64 -> 32). Knowing the model would already struggle with overfitting due to the small sample size, it would be counter-intuitive to use many layers. 



\subsection{Results}

\subsubsection{Development and training time}

\placeholder{Talk about dev and training time}

\subsubsection{Numerical metrics}

\placeholder{Talk about numerical metrics (F1, roc auc, etc)}

\subsubsection{Confusion matrix}

\placeholder{Talk about confusion matrix}

\subsubsection{Predicted results on test set}

\placeholder{Talk about predicted results on test set}


\subsection{Conclusion}

\subsubsection{Lessons learned}

\placeholder{What have you learned?}


% =============================================================================
\end{document}
